% Part: incompleteness
% Chapter: arithmetization-syntax
% Section: coding-terms

\documentclass[../../../include/open-logic-section]{subfiles}

\begin{document}

\olfileid[pt]{inc}{art}{trm}
\olsection{Codificação de termos}

\begin{explain}
Um termo é simplesmente certo tipo de sequência de símbolos: ele é
construído indutivamente a partir de constantes e variáveis, de acordo
com as regras de formação de termos. Como sequências de símbolos podem
ser codificadas como números---mediante um esquema de codificação dos
símbolos e uma maneira de codificar sequências de números---não é
difícil atribuir números de G\"odel aos termos. O desafio é mostrar que
a propriedade que um número possui quando é o número de G\"odel de um
termo corretamente formado é computável ou, na verdade, recursiva
primitiva.
\end{explain}

!!^{variable}s e !!{constant}s são os termos mais simples, e é fácil
testar se $x$ é o número de G\"odel de um termo desse tipo:
$\fn{Var}(x)$ vale se $x$ é $\Gn{\Obj v_i}$ para algum~$i$. Em outras
palavras, $x$~é uma sequência de comprimento~$1$ e seu único elemento
$(x)_0$ é o código de alguma !!{variable}~$\Obj v_i$, O.s.\@, $x$ é
$\tuple{\tuple{1, i}}$ para algum~$i$. De modo semelhante,
$\fn{Const}(x)$ vale se $x$ é $\Gn{\Obj c_i}$ para algum~$i$. Ambas as
relações são recursivas primitivas, pois, se tal $i$ existe, ele deve
ser $< x$:
\begin{align*}
  \fn{Var}(x) & \defiff \bexists{i<x}{x = \tuple{\tuple{1, i}}}\\
  \fn{Const}(x) & \defiff \bexists{i<x}{x = \tuple{\tuple{2, i}}}
\end{align*}

\begin{prop}
\ollabel{prop:term-primrec}
As relações $\fn{Term}(x)$ e $\fn{ClTerm}(x)$, que valem sse $x$ é,
respectivamente, o número de G\"odel de um termo ou de um termo fechado,
são recursivas primitivas.
\end{prop}

\begin{proof}
Uma sequência de símbolos~$s$ é um termo sse existe uma sequência de
termos~$s_0$, \dots, $s_{k-1} = s$ que registra como o termo~$s$ foi
formado a partir de !!{constant}s e !!{variable}s segundo as regras de
formação de termos. Para expressar que uma suposta sequência de
formação segue essas regras, deve ocorrer, para cada $i < k$, uma das
seguintes alternativas:
\begin{enumerate}
\item $s_i$ é !!a{variable}~$\Obj v_j$, ou
\item $s_i$ é !!a{constant}~$\Obj c_j$, ou
\item $s_i$ é construído a partir de $n$ termos $t_1$, \dots, $t_n$
  que ocorrem antes da posição~$i$, usando !!a{function} de $n$
  lugares~$\Obj f^n_j$.
\end{enumerate}
Para mostrar que a relação correspondente sobre números de G\"odel é
recursiva primitiva, precisamos expressar essa condição recursivamente
de modo primitivo, O.s.\@, usando funções e relações recursivas
primitivas e quantificação limitada.

Suponha que $y$ seja o número que codifica a sequência $s_0$, \dots,
$s_{k-1}$, O.s.\@, $y = \tuple{\Gn{s_0}, \dots, \Gn{s_{k-1}}}$. Ele
codifica uma sequência de formação para o termo com número de
G\"odel~$x$ sse, para todo $i < k$:
\begin{enumerate}
\item $\fn{Var}((y)_i)$, ou
\item $\fn{Const}((y)_i)$, ou
\item existem um $n$ e um número~$z = \tuple{z_1, \dots, z_n}$ tais
  que cada $z_l$ é igual a algum $(y)_{i'}$ para $i' < i$ e
\[
(y)_i = \Gn{\Obj f^n_j(} \concat \fn{flatten}(z) \concat \Gn{)},
\]
\end{enumerate}
e, além disso, $(y)_{k-1} = x$. (A função $\fn{flatten}(z)$ transforma
a sequência $\tuple{\Gn{t_1}, \dots, \Gn{t_n}}$ em $\Gn{t_1, \dots,
  t_n}$ e é recursiva primitiva.)

Os índices $j$, $n$, os números de G\"odel $z_l$ dos termos $t_l$ e o
código~$z$ da sequência~$\tuple{z_1, \dots, z_n}$, em (3), são todos
menores que~$y$. Podemos substituir o $k$ acima por $\len{y}$. Assim,
podemos expressar ``$y$ é o código de uma sequência de formação do
termo com número de G\"odel~$x$'' de uma maneira que mostra que essa
relação é recursiva primitiva.

Resta apenas nos convencermos de que existe um limite recursivo
primitivo para~$y$. Mas, se $x$ é o número de G\"odel de um termo, ele
deve ter uma sequência de formação com no máximo $\len{x}$ termos
(pois cada termo da sequência de formação de~$s$ deve começar em alguma
posição de~$s$, e dois subtermos não podem começar na mesma posição).
O número de G\"odel de cada subtermo de~$s$ é, naturalmente, $\le x$.
Logo, sempre existe uma sequência de formação com código
$\le p_{k-1}^{k(x+1)}$, em que $k=\len{x}$.

Para $\fn{ClTerm}$, basta omitir a cláusula referente às !!{variable}s.
\end{proof}

\begin{prob}
Mostre que a função $\fn{flatten}(z)$, que transforma a sequência
$\tuple{\Gn{t_1}, \dots, \Gn{t_n}}$ em $\Gn{t_1, \dots, t_n}$, é
recursiva primitiva.
\end{prob}

\begin{prop}\ollabel{prop:num-primrec}
  A função $\fn{num}(n) = \Gn{\num{n}}$ é recursiva primitiva.
\end{prop}

\begin{proof}
  Definimos $\fn{num}(n)$ por recursão primitiva:
  \begin{align*}
    \fn{num}(0) & = \Gn{\Obj 0}\\
    \fn{num}(n+1) & = \Gn{\prime(} \concat \fn{num}(n) \concat \Gn{)}.
  \end{align*}
\end{proof}

\end{document}
